% Source: ACM_Templates.pdf, PDF pp. 21-27.
% Corrections and validation: docs/validation.md.
\subsection{矩阵快速幂与广义矩阵乘法}
\label{algo:matrix}
\subsubsection*{矩阵怎样表示状态转移}
本节统一使用\textbf{行向量}：$v_{t+1}=v_tM$。矩阵的行表示旧状态、列表示新状态，所以
\[
    (vM)_j=\sum_i v_iM_{ij},\qquad
    (AB)_{ij}=\sum_k A_{ik}B_{kj}.
\]
先执行 $A$、再执行 $B$，对应 $AB$；矩阵乘法通常不可交换。相同转移连续执行 $k$ 次就是 $v_0M^k$，利用乘法结合律做二进制快速幂。$d$ 阶矩阵每次乘法 $O(d^3)$，快速幂 $O(d^3\log(k+1))$，空间 $O(d^2)$。

以 $f_t=f_{t-1}+f_{t-3}$ 为例，取 $v_t=[f_t,f_{t-1},f_{t-2}]$，则
\[
    M=\begin{pmatrix}1&1&0\\0&0&1\\1&0&0\end{pmatrix},
    \qquad v_n=v_3M^{n-3}\quad(n\geq3).
\]
构造时逐列写出新状态依赖哪些旧状态，能减少行、列向量约定混用的错误。常数项可通过附加一个恒为 $1$ 的状态处理。

\subsubsection*{普通模矩阵的使用教程}
\InlineCode{Mat(n,mod)} 创建从 $1$ 开始的 $n$ 阶零矩阵，第三个参数为真时创建单位矩阵。通过 \InlineCode{a[i][j]} 赋值，\InlineCode{pow(k)} 返回非负整数次幂，不修改原对象。参与相乘的维数和模数必须相同，元素预先归一化到 $[0,mod-1]$，且 $1\leq mod\leq2\times10^9$。
\begin{lstlisting}[style=mtycpp,numbers=none]
Mat m(2,1000000007);
m.a[1][1]=1,m.a[1][2]=1,m.a[2][1]=1;
auto r=m.pow(10);
\end{lstlisting}

\needspace{7\baselineskip}
\subsubsection*{模板代码}
\CodeFile{code/number-theory/matrix.cpp}

\subsubsection*{广义乘法：用最优值替代方案数}
只要运算满足需要的结合律、分配律和单位元，就能将矩阵中的求和与相乘替换为其他运算。常见模型为
\begin{center}
\small
\begin{tabular}{@{}llll@{}}
\toprule
模型 & 合并备选 & 连接两段 & 零元 / 单位元 \\
\midrule
计数 & 加法 & 乘法 & $0$ / $1$ \\
最短路 & $\min$ & 加法 & $+\infty$ / $0$ \\
最大收益 & $\max$ & 加法 & $-\infty$ / $0$ \\
可达性 & 逻辑或 & 逻辑与 & 假 / 真 \\
\bottomrule
\end{tabular}
\end{center}
在 max-plus 模型中，$M_{ij}$ 是一步从 $i$ 到 $j$ 的最大收益，不可达为 $-\infty$。$M^k$ 表示\textbf{恰好}走 $k$ 步的最大收益；若要求至多 $k$ 步，必须另行建模。单位矩阵对角线为 $0$，其余为 $-\infty$。

\subsubsection*{max-plus 模板的使用教程}
\InlineCode{MaxMat(n)} 创建所有转移都不可达的矩阵，第二个参数为真时创建单位矩阵。普通赋值、乘法、幂接口与上面的 \InlineCode{Mat} 一致。\InlineCode{INF=-4e18} 是不可达哨兵，代码在相加前跳过它；所有真实路径收益需严格位于 $(-4\times10^{18},4\times10^{18})$，避免与哨兵冲突并保证中间加法安全。
\begin{lstlisting}[style=mtycpp,numbers=none]
MaxMat m(2);
m.a[1][1]=0,m.a[1][2]=5,m.a[2][2]=2;
auto r=m.pow(3);
\end{lstlisting}

\needspace{7\baselineskip}
\subsubsection*{模板代码}
\CodeFile{code/number-theory/max_matrix.cpp}

\subsubsection*{应用：最大子段和的转移矩阵}
使用三种状态：尚未开始、正在选择子段、已经结束。对当前值 $x$，行向量转移为
\[
    T(x)=\begin{pmatrix}
        0&x&x\\
        -\infty&x&x\\
        -\infty&-\infty&0
    \end{pmatrix}.
\]
从 $[0,-\infty,-\infty]$ 出发，读取状态 $3$ 即为非空最大子段和：每个被选的 $x$ 都累加一次，进入结束状态后只能跳过后续元素。矩阵的 $(1,3)$ 项允许长度为 $1$ 的子段直接开始并结束。

区间 $[l,r]$ 对应有序乘积 $T(a_l)\cdots T(a_r)$。静态查询可用保持顺序、互不重叠的倍增块拼接；带修改时在线段树中维护左右儿子的乘积。此乘法不满足幂等性，不能像普通最值 ST 表一样用两个重叠块回答查询。
